%% Part 2.2 depth entries and lookup tables.
%% Register coverage: D-P3-003, D-P3-004, D-P3-006, D-P3-014;
%% apparatus A-035 and A-036.  D-P3-005 and D-P3-007 are pointers only
%% and deliberately remain open at their Part-3 and Part-2.4 homes.

\chapter{First-Order Reference Entries}
\label{sec:p02-first-order-depth}

\noindent\textit{Reader note.}  This chapter is an editorial reference layer
after the exact inherited first-order sequence.  Its entries and tables retain
their own citations and do not replace, rename, or change the proof status of
any inherited statement.

This chapter collects semantic reductions and metatheorem lookups that do not
belong to a single calculus presentation.  Herbrand's theorem is stated in its
validity form; its clause-set form remains with resolution.  Interpolation,
definability, and the quantified subformula property are citation-only entries:
no proof status is inferred from their inclusion here.

\section{Herbrand universe, structures, and finite expansions}
\label{sec:p02-herbrand-validity}

Fix an equality-free first-order signature \(\Lang\) with at least one ground
term.  If \(\Lang\) has none, first add one fresh constant \(c\).  A sentence
of the old signature is valid exactly when it is valid in that fresh-constant
expansion, because every expanded structure has an old-signature reduct and
every nonempty old-signature structure admits an interpretation of \(c\).

\begin{hbdefinition}{def:p02-herbrand-universe-base}{Herbrand universe and Herbrand base}
The \emph{Herbrand universe} of \(\Lang\) is the set of all ground
\(\Lang\)-terms.  The \emph{Herbrand base} is the set of all ground atoms
\[
  R(t_1,\ldots,t_n),
\]
where \(R\) is an \(n\)-ary relation symbol of \(\Lang\) and every \(t_i\)
belongs to the Herbrand universe.  Nullary relation symbols contribute their
own atoms.  Logical equality is excluded from this base and from the theorem
below; it is not treated as an arbitrary binary predicate.
\index{Herbrand universe}
\index{Herbrand base}
\end{hbdefinition}

\begin{hbdefinition}{def:p02-herbrand-structure}{Herbrand structure}
For a set \(B\) of members of the Herbrand base, the corresponding
\emph{Herbrand structure} \(\mathcal M_B\) has the Herbrand universe as its
domain, interprets each constant \(d\) by the ground term \(d\), and interprets
each \(n\)-ary function symbol by term formation:
\[
 f^{\mathcal M_B}(t_1,\ldots,t_n)
   \defeq f(t_1,\ldots,t_n).
\]
Its relations are fixed by
\[
 \mathcal M_B\entails R(t_1,\ldots,t_n)
 \quad\Longleftrightarrow\quad
 R(t_1,\ldots,t_n)\in B.
\]
Thus a Herbrand structure varies only in which ground relation atoms it makes
true; constants and functions always have their syntactic interpretations.
\index{Herbrand structure}
\end{hbdefinition}

\begin{hbdefinition}{def:p02-herbrand-expansion}{Finite Herbrand expansion}
Let
\[
 \sigma\defeq\exists x_1\cdots\exists x_n\,\varphi(x_1,\ldots,x_n)
\]
be a sentence with quantifier-free matrix \(\varphi\), where \(n\geq0\).  Given
finitely many \(n\)-tuples \(\bar t_1,\ldots,\bar t_m\) of ground terms, with
\(m\geq1\), the associated \emph{finite Herbrand expansion} is
\[
 \varphi(\bar t_1)\lor\cdots\lor\varphi(\bar t_m).
\]
Each substitution is the total capture-avoiding substitution fixed in the
syntax chapter.  For propositional evaluation, each distinct ground atom is
regarded as one propositional variable.  When \(n=0\), the tuple is empty and
the quantifier-free case is the trivial one-member expansion \(\varphi\) (take
\(m=1\)).
\index{Herbrand expansion}
\end{hbdefinition}

\begin{hbtheorem}{thm:p02-herbrand-validity}{cited=AfshariEnqvistLeigh2025}{Herbrand's theorem, validity form}
Let \(\Lang\) be equality-free and contain a ground term, and let
\(\sigma=\exists\bar x\,\varphi(\bar x)\) be a prenex existential \(\Lang\)-sentence
with quantifier-free matrix.  Then \(\sigma\) is first-order valid iff some
finite Herbrand expansion of \(\sigma\) is propositionally valid.  If the
original signature has no ground term, the same equivalence holds after adding
one fresh constant as above.
\index{Herbrand's theorem!validity form}
\end{hbtheorem}

Both directions are proved in
Theorem~\ref{thm:p02-w02-p02-028-existential-validity}, and
Theorem~\ref{thm:p02-w02-p02-028-general-validity} extends the equivalence to
arbitrary equality-free sentences through the Herbrand form.
The theorem does not say that validity supplies one uniformly valid ground
instance: a finite disjunction of instances may be necessary.  Nor does it
identify three different constructions.  Theorem~\ref{thm:herbrand-clauses}
is the dual finite-ground-instance criterion for unsatisfiable universal clause
sets.  The term structure in Section~\ref{sec:fol-metatheory-02} instead quotients all
terms by provable equality in a maximal consistent witness theory.  A Herbrand
structure is an unquotiented ground-term structure and the validity theorem is
equality-free; neither is the Henkin term model.

\section{Interpolation and implicit definability}
\label{sec:p02-craig-beth}

\begin{hbtheorem}{thm:p02-craig-interpolation}{cited=Craig1957}{Craig interpolation, sentence form}
Let \(\varphi\) and \(\psi\) be first-order sentences, possibly in different
signatures, and suppose \(\entails\varphi\to\psi\).  Then there is a sentence
\(\chi\) such that
\[
 \entails\varphi\to\chi,
 \qquad
 \entails\chi\to\psi,
\]
and every nonlogical constant, function symbol, and relation symbol occurring
in \(\chi\) occurs in both \(\varphi\) and \(\psi\).  When equality is logical,
it is exempt from the shared-vocabulary condition.
\index{Craig interpolation theorem}
\end{hbtheorem}

\begin{hbdefinition}{def:p02-beth-definability}{Implicit and explicit predicate definability}
Let \(\Lang(P)\) extend \(\Lang\) by an \(n\)-ary predicate \(P\),
let \(T(P)\) be a \(\Lang(P)\)-theory, let
\(x_1,\ldots,x_n\) be pairwise distinct, and let \(P'\) be fresh.
Write \(T(P')\) for the uniform replacement of \(P\) by \(P'\).  The theory
\emph{implicitly defines} \(P\) over \(\Lang\) iff
\[
 T(P)\cup T(P')\entails
 \forall x_1\cdots\forall x_n
 \bigl(P(x_1,\ldots,x_n)\leftrightarrow
       P'(x_1,\ldots,x_n)\bigr).
\]
It \emph{explicitly defines} \(P\) over \(\Lang\) iff some
\(\Lang\)-formula \(\chi\), whose free variables are among
\(x_1,\ldots,x_n\), satisfies
\[
 T(P)\entails
 \forall x_1\cdots\forall x_n
 \bigl(P(x_1,\ldots,x_n)\leftrightarrow
       \chi(x_1,\ldots,x_n)\bigr).
\]
These are metalinguistic properties of a theory and its distinguished
predicate, not new object-language connectives.
\index{implicit definability}
\index{explicit definability}
\end{hbdefinition}

\begin{hbtheorem}{thm:p02-beth-definability}{cited=Beth1953}{Beth definability theorem}
For every first-order theory \(T(P)\) as in
Definition~\ref{def:p02-beth-definability}, \(T(P)\) implicitly defines \(P\)
over \(\Lang\) iff it explicitly defines \(P\) over \(\Lang\).
\index{Beth definability theorem}
\end{hbtheorem}

Complete proofs of both theorems are given elsewhere in the book.  For
interpolation, Theorem~\ref{thm:craig-interpolation} is proved in
Annex~\ref{annex:p7.interp.3} by Maehara's partition method on cut-free
derivations, and Theorem~\ref{thm:p03-craig-robinson-beth}
(Annex~\ref{annex:p3.abstract.5}) derives it from Robinson joint consistency,
Theorem~\ref{thm:p03-robinson-joint-consistency}
(Annex~\ref{annex:p3.abstract.1}).  For definability,
Theorem~\ref{thm:beth} is proved in Annex~\ref{annex:p7.interp.4} and
Theorem~\ref{thm:p03-beth-definability} in Annex~\ref{annex:p3.abstract.2}.
The formula form of interpolation, with its free-variable condition, is
Theorem~\ref{thm:p02-w02-p02-037-craig-formula} in
\cref{sec:p02-w02-p02-037-craig}.  \Cref{sec:p02-w02-p02-037-beth} shows
that Definition~\ref{def:p02-beth-definability} agrees with the model
formulation used by those proofs and derives
Theorem~\ref{thm:p02-beth-definability} as
Corollary~\ref{thm:p02-w02-p02-037-beth}.  Lyndon interpolation is not
claimed here.

\section{Decidable fragments: exact scopes}
\label{sec:p02-decidable-fragments}

Decision and complexity claims are sensitive to vocabulary, equality,
function symbols, model class, and whether the input problem is satisfiability
or validity.  Table~\ref{tab:p02-decidable-fragments} fixes those parameters
row by row.  Its catalog is a synthesis of the cited results, not a claim that
the fragment names retain the displayed bounds under arbitrary extensions.

\begin{hbdefinition}{def:p02-exponential-time-classes}{Exponential-time classes used in this table}
For input length \(n\), let
\[
\begin{aligned}
 \mathsf{EXPTIME}&\defeq\bigcup_{k\geq1}\mathsf{DTIME}(2^{n^k}),&
 \mathsf{NEXPTIME}&\defeq\bigcup_{k\geq1}\mathsf{NTIME}(2^{n^k}),\\
 \mathsf{coNEXPTIME}&\defeq
   \{L:\overline L\in\mathsf{NEXPTIME}\},&
 2\mathsf{EXPTIME}&\defeq
   \bigcup_{k\geq1}\mathsf{DTIME}(2^{2^{n^k}}).
\end{aligned}
\]
Here \(\mathsf{DTIME}\) and \(\mathsf{NTIME}\) use deterministic and
nondeterministic multitape Turing machines under the standard finite encoding
of formulas.  These names are needed because the table distinguishes
determinism, nondeterminism, complementation, and one versus two exponential
levels.  They denote decision-problem complexity classes, not logical systems,
and are used only with the reductions and input languages stated row by row.
\index{EXPTIME}\index{NEXPTIME}\index{2EXPTIME}
\end{hbdefinition}

\begin{landscape}
\begingroup\scriptsize\setlength{\tabcolsep}{3pt}
\begin{longtable}{@{}p{.90in}p{2.20in}p{2.10in}p{2.00in}@{}}
\caption{Decidable fragments of first-order logic.}
\label{tab:p02-decidable-fragments}\\
\toprule
Fragment & Exact input language and semantics & Model and decision complexity & Boundary and source \\ \midrule
\endfirsthead
\caption[]{Decidable fragments of first-order logic (continued)}\\
\toprule
Fragment & Exact input language and semantics & Model and decision complexity & Boundary and source \\ \midrule
\endhead
\midrule\multicolumn{4}{r}{\emph{continued on next page}}\\
\endfoot
\endlastfoot
Monadic FOL
\index{monadic first-order logic} &
Closed formulas over unary predicate letters, with nonempty-domain semantics;
Lewis's language has neither identity nor function signs, and no constants are
imported as zero-ary functions. &
Satisfiability is decidable and \(\mathsf{NEXPTIME}\)-complete; more sharply,
Lewis gives
matching nondeterministic bounds of the form \(c^{n/\log n}\), with different
constants.  Validity is \(\mathsf{coNEXPTIME}\)-complete by complementation
\cite{Lewis1980}.  Without equality every satisfiable sentence with \(m\)
predicates has a model with at most \(2^m\) elements, so the row has the
finite-model property; the colour-set procedure and the matching lower bound are in
\cref{sec:p02-w02-p02-035-monadic}. &
Equality, constants, and function symbols are outside this cited complexity
claim; adding them is not a notational variant.
Undecidable boundary: a single binary predicate symbol already yields a
conservative reduction class \cite[thm.~3.1.23]{BorgerGradelGurevich1997}. \\
\addlinespace
Two-variable logic \(\mathrm{FO}^{2}\)
\index{two-variable logic} &
Relational first-order sentences using at most two distinct variable symbols,
reused under quantifiers; the theorem holds with or without logical equality.
No function symbols are present.  Individual constants may be added; they are
kings, and the model and complexity bounds are unchanged
\cite[p.~11]{GradelKolaitisVardi1997}. &
Finite-model property; every satisfiable sentence has a model of exponential
size.  Finite and unrestricted satisfiability coincide.  Satisfiability, and
hence finite satisfiability, is
\(\mathsf{NEXPTIME}\)-complete \cite{GradelKolaitisVardi1997}.
Construction and algorithm:
\cref{sec:p02-w02-p02-035-fo2,sec:p09-presentation-n-c-006}. &
The three-variable relational fragment already has an undecidable
satisfiability problem; the two-variable bound is therefore load-bearing
\cite{GradelKolaitisVardi1997}. \\
\addlinespace
Guarded fragment\newline\(\mathrm{GF}\)
\index{guarded fragment} &
Relational first-order formulas in the cited presentation, with no function
symbols.  Logical equality may occur as an atom but may not itself serve as a
guard.  Every quantified tuple is relativized by a non-equality relation atom
containing all variables free in the quantified subformula.
Individual constants may occur in the input; a polynomial-time translation
removes them and preserves satisfiability over every universe
\cite[lem.~2.3]{Gradel1999} (\cref{thm:p02-w02-p02-035-gf-constants}).
That translation raises arity and variable count; the bounded-variable and
bounded-arity \(\mathsf{EXPTIME}\) bounds hold with constants by a different
reduction (\cref{thm:p02-w02-p02-035-gf-bounded-constants}). &
Tree-model property and finite-model property; satisfiability is decidable and
general satisfiability is \(2\mathsf{EXPTIME}\)-complete.  With a fixed bound
on the number of variables or relation arity it is
\(\mathsf{EXPTIME}\)-complete \cite{Gradel1999}.
Tree-model and finite-model constructions, type-elimination procedure and
lower bounds: \cref{sec:p02-w02-p02-035-gf}; source presentation:
\cref{sec:p09-presentation-n-c-007}. &
An atom beside a quantifier is not a guard unless it covers all relevant
variables.  Loosely guarded and fixed-point extensions are different rows.
Undecidable boundary: adding functionality statements, or transitivity
statements, to the three-variable guarded fragment gives conservative
reduction classes \cite[thms.~5.8 and 5.10]{Gradel1999}. \\
\addlinespace
Bernays--Sch\"onfinkel\newline(\(\mathrm{EPR}\))
\index{Bernays--Schonfinkel class@Bernays--Sch\"onfinkel class}
\index{effectively propositional logic} &
Closed, function-free relational sentences in prenex
\(\exists^{*}\forall^{*}\) form with quantifier-free matrix.  The cited
complexity convention omits identity; Skolemizing the leading existentials
introduces constants only. &
Finite-model property; satisfiability is decidable by finite grounding and is
\(\mathsf{NEXPTIME}\)-complete \cite{Lewis1980}.
Grounding procedure, correctness and bounds:
\cref{sec:p02-w02-p02-035-bsr}. &
The reversed prefix, logical equality, and positive-arity function symbols are
not included in this row's theorem.
Undecidable boundary: the equality-free prefix classes
\(\forall^{3}\exists\) and \(\forall\exists\forall\) are reduction classes
\cite[p.~311]{BorgerGradelGurevich1997}. \\
\addlinespace
Presburger arithmetic
\index{Presburger arithmetic} &
Sentences of the complete first-order theory of
\((\N,0,1,+,<)\), with logical equality; multiplication of two variables is
absent.  Integer presentations require an explicit standard intertranslation. &
Decidable by quantifier elimination \cite{Presburger1929}.  Classically, a
nondeterministic doubly-exponential lower bound
\cite{FischerRabin1974} and a deterministic triply-exponential upper bound
\cite{Oppen1978}; this row does not assert a tight complexity class. &
Adding multiplication of two variables yields the full first-order theory of
\((\N,0,1,+,\times,<)\), whose decision problem is undecidable
\cite{DavisPutnamRobinson1961,Matiyasevich1970}; changing numeral encoding or the domain
convention changes the quoted complexity problem. \\
\bottomrule
\end{longtable}
\endgroup
\end{landscape}

The boundary column names an undecidable neighbouring class for each row---one
binary predicate beyond the monadic class, three variables beyond
\(\mathrm{FO}^{2}\), functionality or transitivity statements beyond the
guarded fragment, the prefixes \(\forall^{3}\exists\) and \(\forall\exists\forall\)
beyond the Bernays--Sch\"onfinkel class, and variable multiplication beyond
Presburger arithmetic---and distinguishes these from features merely excluded
from a cited scope.
Dropping all restrictions reaches unrestricted first-order validity, which is
undecidable \cite{Church1936,Turing1937}; satisfiability is therefore
undecidable as well, by negation.  No row makes the generally false assertion
that every excluded feature, by itself, destroys decidability.
Part~14 lists the same fragments, with their semantic ranges, in
Table~\ref{tab:p14-dc-finite}; the decision procedures for the monadic,
two-variable, guarded and Bernays--Sch\"onfinkel rows are written out in
\cref{sec:p02-w02-p02-035-decision-procedures}.

\section{The quantified subformula property}
\label{sec:p02-fol-subformula-property}

\begin{hbdefinition}{def:p02-fol-subformula-instance}{First-order subformula instance}
The relation ``\(\psi\) is a \emph{first-order subformula instance} of
\(\chi\)'' is the least metalinguistic relation on formulas satisfying:
\begin{enumerate}
  \item \(\chi\) is a first-order subformula instance of itself;
  \item whenever an admitted formula has an immediate propositional
        subformula, that immediate subformula is admitted; and
  \item whenever \(Qx\,\varphi\) is admitted, where
        \(Q\in\{\forall,\exists\}\), every capture-avoiding instance
        \(\varphi[t/x]\) is admitted for every term \(t\) in the language of
        the derivation.
\end{enumerate}
The clauses reapply to newly admitted formulas.  This extension of literal
syntactic subformulahood is necessary because \(\forall L\) and \(\exists R\)
use arbitrary term instances, while \(\forall R\) and \(\exists L\) use fresh
eigenvariables, which are variable terms.  It is a relation among formula
occurrences in proof syntax, not formula equivalence, and it differs from the
propositional operation \(\operatorname{Sub}\) exactly by clause~3.
\index{subformula!first-order instance}
\index{eigenvariable!subformula property}
\end{hbdefinition}

\begin{hbtheorem}{thm:p02-fol-subformula-property}{proved=annex:p02-w02b-p02-020-host-property}{Subformula property for cut-free first-order \(\mathsf{LK}\)}
In every derivation using neither Cut nor Mix on the equality-free registered
sheet \hyperref[sheet:lk]{\texttt{sheet:lk}}, every formula occurrence is a
first-order subformula instance of some formula occurrence in the endsequent.
In particular, the claim includes the term instances introduced by
\(\forall L\) and \(\exists R\), and the fresh eigenvariable instances
introduced by \(\forall R\) and \(\exists L\).  It is not the false claim that
every proof formula is a literal string-subformula of the endsequent, and it
does not apply to a derivation still containing Cut or Mix.
\end{hbtheorem}

\section{Many-sorted destination}
\label{sec:p02-first-order-depth-pointers}

The registered home for many-sorted first-order syntax, nonempty carrier
semantics, and the translation to one-sorted logic is
Section~\ref{sec:p02-many-sorted}.  Its guard translation supplies unary sort
predicates, relativized quantifiers, nonemptiness axioms, and symbol-typing
axioms; it is not the erasure of sort annotations \cite{Manzano1996}.  The
exact reduction is Theorem~\ref{thm:p02-many-sorted-translation}.

\section{First-order metatheorem summary}
\label{sec:p02-fol-metatheorem-summary}

Table~\ref{tab:p02-fol-metatheorems} distinguishes object-level hypotheses from
the set-theoretic assumptions under which this handbook invokes each result.
For a citation-only row, ``ZFC lookup'' means exactly that the result is used in
ordinary ZFC metatheory, so AC is available; it neither claims that AC is
necessary nor asserts a proof over a weaker base theory.

For the Lindstr\"om row, a \emph{regular abstract logic}
\index{regular abstract logic} has this fixed metalinguistic meaning in the
relational presentation.  It assigns to each relational signature \(\tau\) a
class of sentences and an isomorphism-invariant satisfaction relation.  Each
sentence has a finite occurrence signature, and satisfaction depends only on
the reduct to that signature; arity-preserving renamings transport both
sentences and satisfaction.  The logic contains relational first-order logic,
is closed under negation and conjunction, and has the following two exact
operations: abstraction of a fresh constant has the usual first-order
quantifier semantics, and for a fresh unary predicate \(U\), relativization
\(\varphi^{U}\) is true in \(\mathcal M\) exactly when \(\varphi\) is true in
the induced \(\tau\)-structure on the nonempty set \(U^{\mathcal M}\).  This
package is a condition on a logic-and-satisfaction assignment, not an
object-language connective and not the proof-theoretic fragment also called
regular logic.  Its only use here is to make the characterization row
unambiguous; see \cite[Chapter~XIII, \S1]{EbbinghausFlumThomas1994}.

\begin{landscape}
\begingroup\scriptsize\setlength{\tabcolsep}{3pt}
\begin{longtable}{@{}p{1.15in}p{2.15in}p{2.25in}p{1.65in}@{}}
\caption{First-order metatheorem summary.}
\label{tab:p02-fol-metatheorems}\\
\toprule
Result & Exact hypotheses & Lookup conclusion & Choice, evidence, and locator \\ \midrule
\endfirsthead
\caption[]{First-order metatheorem summary (continued)}\\
\toprule
Result & Exact hypotheses & Lookup conclusion & Choice, evidence, and locator \\ \midrule
\endhead
\midrule\multicolumn{4}{r}{\emph{continued on next page}}\\
\endfoot
\endlastfoot
Natural-deduction soundness &
Arbitrary signature; arbitrary \(\Gamma\cup\{\varphi\}\); nonempty set-based
structures; logical equality is identity. &
\(\Gamma\derives_{\mathsf{ND}_{\mathsf{FOL}}}\varphi\) implies
\(\Gamma\entails\varphi\). &
No choice used; complete inherited proof,
Theorem~\ref{thm:fol-soundness}. \\
\addlinespace
Strong completeness &
Arbitrary signature and formula set under AC; alternatively a finite or
countable signature with no choice assumption. &
\(\Gamma\entails\varphi\) implies
\(\Gamma\derives_{\mathsf{ND}_{\mathsf{FOL}}}\varphi\). &
Proved: AC in the arbitrary-signature construction; choice-free in the
countable construction.  Theorem~\ref{thm:fol-completeness} and
Corollary~\ref{cor:fol-completeness-countable}. \\
\addlinespace
Compactness &
Under AC, \(T\) is any set of \(\Lang\)-sentences and \(\sigma\) is a
\(\Lang\)-sentence. &
\(T\) has a model iff every finite subset does; if \(T\entails\sigma\), some
finite \(T_0\subseteq T\) entails \(\sigma\). &
AC in this handbook's derivation; complete inherited proof,
Theorem~\ref{thm:fol-compactness}. \\
\addlinespace
Downward L\"owenheim--Skolem &
Under AC, \(\mathcal M\) is a structure and \(A\subseteq\dom\mathcal M\). &
Some \(\mathcal N\preccurlyeq\mathcal M\) contains \(A\) and has size at most
\(\max(\abs{A},\abs{\Lang},\aleph_0)\). &
AC in the stated theorem; complete inherited proof,
Theorem~\ref{thm:downward-ls}. \\
\addlinespace
Upward L\"owenheim--Skolem &
Under AC, \(T\) is an \(\Lang\)-theory with an infinite model and
\(\kappa\) is a cardinal with \(\kappa\geq\max(\abs{\Lang},\aleph_0)\). &
\(T\) has a model of cardinality exactly \(\kappa\). &
AC in the stated theorem; complete inherited proof,
Theorem~\ref{thm:upward-ls}. \\
\addlinespace
Herbrand validity &
Equality-free prenex existential sentence with quantifier-free matrix;
signature has a ground term, or is expanded by one fresh constant. &
Validity iff some finite ground expansion is propositionally valid. &
No choice principle in the proof given here.  Both directions are proved as
Theorem~\ref{thm:p02-w02-p02-028-existential-validity}
(Annex~\ref{annex:p02-w02-p02-028-existential-validity}): left to right, the
proof passes to the finite signature of \(\sigma\) and applies
Lemma~\ref{thm:p02-w02-p02-028-herbrand-model}, whose one infinitary step is
Theorem~\ref{thm:prop-compactness}, obtained from
Theorem~\ref{thm:prop-completeness} and Lemma~\ref{lem:prop-lindenbaum},
which recurses along a fixed enumeration of the countable set
\(\mathsf{Fm}_{\mathsf{CPL}}\); right to left uses only ground evaluation.
Statement Theorem~\ref{thm:p02-herbrand-validity};
\cite{AfshariEnqvistLeigh2025}.  The extension to arbitrary equality-free
sentences, Theorem~\ref{thm:p02-w02-p02-028-general-validity}, uses AC right
to left, through Theorem~\ref{thm:skolem-equisatisfiable}(1). \\
\addlinespace
Craig interpolation &
\(\varphi,\psi\) are first-order sentences, possibly in different
signatures, and \(\entails\varphi\to\psi\); every nonlogical constant,
function symbol and relation symbol of the interpolant occurs in both
\(\varphi\) and \(\psi\), and logical equality is exempt from this
shared-vocabulary condition. &
There is \(\chi\) with \(\entails\varphi\to\chi\) and
\(\entails\chi\to\psi\). &
ZFC lookup: AC is available and no weaker-base claim is made.  Statement
Theorem~\ref{thm:p02-craig-interpolation}; \cite{Craig1957}; formula form
with free-variable condition
Theorem~\ref{thm:p02-w02-p02-037-craig-formula}; proofs
Theorem~\ref{thm:craig-interpolation} (Annex~\ref{annex:p7.interp.3}) and
Theorem~\ref{thm:p03-craig-robinson-beth} (Annex~\ref{annex:p3.abstract.5}). \\
\addlinespace
Beth definability &
An \(\Lang(P)\)-theory; implicit uniqueness is expressed by a renamed
copy \(P'\) over the shared \(\Lang\)-reduct. &
Implicit definability of \(P\) iff a \(\Lang\)-formula explicitly
defines \(P\) modulo the theory. &
ZFC lookup: AC is available and no weaker-base claim is made.  Statement
Theorem~\ref{thm:p02-beth-definability}; \cite{Beth1953}; proved in this
formulation as Corollary~\ref{thm:p02-w02-p02-037-beth}, from
Theorem~\ref{thm:beth} (Annex~\ref{annex:p7.interp.4}) or
Theorem~\ref{thm:p03-beth-definability} (Annex~\ref{annex:p3.abstract.2});
Robinson route Theorem~\ref{thm:p03-robinson-joint-consistency}
(Annex~\ref{annex:p3.abstract.1}) and
Theorem~\ref{thm:p03-craig-robinson-beth} (Annex~\ref{annex:p3.abstract.5}). \\
\addlinespace
Lindstr\"om characterization &
A regular abstract logic in the exact preceding sense.  Every countable set of
its sentences whose finite subsets are satisfiable is satisfiable, and every
satisfiable sentence has a countable model. &
Every sentence of the abstract logic is equivalent, over its occurrence
signature, to a relational first-order sentence; the converse holds because
the logic extends first-order logic. &
ZFC lookup: AC is available and no weaker-base claim is made; cited statement
from \cite{Lindstrom1969,EbbinghausFlumThomas1994}; the public destination is
Theorem~\ref{thm:p03-lindstrom}. \\
\bottomrule
\end{longtable}
\endgroup
\end{landscape}
